% Discovery Tetrad architecture: 4 gates × 2 layers + cascade asymmetry.
% Designed for full text-width in NeurIPS two-column layout.
\begin{figure}[t]
\centering
\begin{tikzpicture}[
  font=\small,
  gate/.style={
    rectangle, rounded corners=2pt, draw=black!60, fill=black!5,
    minimum width=2.2cm, minimum height=0.8cm, align=center
  },
  layerlbl/.style={
    rectangle, draw=black!50, fill=black!8,
    minimum width=1.8cm, minimum height=0.6cm, align=center
  },
  cellbox/.style={
    rectangle, draw=black!40, fill=white,
    minimum width=2.0cm, minimum height=0.55cm, align=center,
    font=\footnotesize
  },
  cascadearrow/.style={
    -{Latex[length=2mm]}, dashed, black!60, line width=0.4pt
  },
  feed/.style={-{Latex[length=2mm]}, black!70, line width=0.5pt},
]

% Gates (top row)
\node[gate] (pl)   at (0,0)   {Plausibility \\ \textit{defeatable?}};
\node[gate] (coh)  at (3,0)   {Coherence \\ \textit{composable?}};
\node[gate] (comp) at (6,0)   {Completeness \\ \textit{template?}};
\node[gate] (inv)  at (9,0)   {Invariance \\ \textit{directional?}};

% Goal layer (left side)
\node[layerlbl] (goal) at (-2.5, -1.7) {\textbf{Goal layer}};
\node[cellbox] (gout) at (-2.5, -2.6) {16 cells};

% Goal layer arrows (independent firing — solid arrows from each gate)
\foreach \g in {pl, coh, comp, inv} {
  \draw[feed] (\g.south) -- ++(0,-0.4) -| (gout.north);
}
\node[font=\scriptsize, black!60] at (3, -1.3) {independent firing};

% Witness layer (right side, lower)
\node[layerlbl] (wit) at (-2.5, -4.8) {\textbf{Witness layer}};
\node[cellbox] (wout) at (-2.5, -5.7) {9 cells};

% Witness layer cascade: Pl/Coh/Comp fire, Inv depends on them
\draw[feed] (pl.south)   -- ++(0,-3.6) -| (wout.north);
\draw[feed] (coh.south)  -- ++(0,-3.7) -| (wout.north);
\draw[feed] (comp.south) -- ++(0,-3.8) -| (wout.north);

% Cascade dependency: Pl/Coh/Comp -> Inv (dashed)
\draw[cascadearrow] (pl.south east)   .. controls +(0.4,-0.6) and +(-1.5,0.3) .. (inv.west);
\draw[cascadearrow] (coh.south east)  .. controls +(0.4,-0.6) and +(-1.0,0.2) .. (inv.west);
\draw[cascadearrow] (comp.south east) .. controls +(0.4,-0.6) and +(-0.5,0.1) .. (inv.west);

\draw[feed] (inv.south) -- ++(0,-3.9) -| (wout.north);

\node[font=\scriptsize, black!60] at (3, -4.5)
  {cascade: Inv fires only if Pl=Coh=Comp=1};

% Cells 8/9 layer-specialization annotation
\node[font=\scriptsize, align=left, black!70] at (12.5, -2.6)
  {bits $(1,1,1,0)$ $\to$ \texttt{supporting\_argument}\\
   bits $(1,1,1,1)$ $\to$ \texttt{fundamental\_result}};
\node[font=\scriptsize, align=left, black!70] at (12.5, -5.7)
  {bits $(1,1,1,0)$ $\to$ \texttt{merely\_true}\\
   bits $(1,1,1,1)$ $\to$ \texttt{established\_result}};

\end{tikzpicture}
\caption{The Discovery Tetrad: four invariant gates applied at two layers.
At the goal layer, all four gates fire independently, producing a 16-cell
verdict taxonomy. At the witness layer, Plausibility, Coherence, and
Completeness fire independently, while Invariance fires only if all three
prerequisites pass --- collapsing the verdict space to nine cells.
The same verdict bits index different cell labels at the two layers because
\textit{transport survival} (witness) and \textit{foundation potential}
(goal) are layer-specific operationalizations of the same abstract gate
question (\textit{directional?}).}
\label{fig:tetrad_arch}
\end{figure}
